% SOURCE_STATEMENT: "What structural mechanism produces the QQ-tilde system for spectral determinants of affine opers, and is there an analogous system on the integrable-model side?"
% Preparation and final self-audit: CONTRIBUTING.md prompts C and D; source/audit artifacts retained in the local physics_statements_2026_10_11 batch.
% !TeX program = lualatex
\documentclass[11pt]{article}
\usepackage[margin=25mm]{geometry}
\usepackage{fontspec}
\setmainfont{Latin Modern Roman}
\usepackage{amsmath,amssymb,mathtools}
\usepackage{unicode-math}
\setmathfont{Latin Modern Math}
\usepackage{xurl,hyperref}
\hypersetup{colorlinks=true,urlcolor=blue}
\setcounter{secnumdepth}{0}
\setlength{\parindent}{0pt}
\setlength{\parskip}{5pt}
\setlength{\emergencystretch}{3em}
\newcommand{\sref}[2]{\hyperlink{source:#1:#2}{[S#2]}}
\newcommand{\cataloguescope}[1]{\paragraph{Recorded target.}#1}
\newcommand{\R}{\mathbb R}
\newcommand{\C}{\mathbb C}
\newcommand{\N}{\mathbb N}
\newcommand{\Z}{\mathbb Z}
\newcommand{\Q}{\mathbb Q}
\newcommand{\F}{\mathbb F}
\newcommand{\eps}{\varepsilon}
\begin{document}
\section{QQ Systems for Spectral Determinants and Integrable Models}\label{statement}

\subsection{Definitions and mathematical statement}

\paragraph{Shared notation.}$\mathbb N=\{1,2,\ldots\}$; the standard number fields have their usual meanings. All additional conventions are stated locally.


\paragraph{Definitions and conventions.}
An oper is a differential-geometric connection with additional flag and transversality structure; affine opers in the source generalize spectral ordinary differential equations. Spectral determinants encode their spectral or connection data. The $Q\widetilde Q$ system is the source\textquotesingle s functional relation among two families of such functions, indexed by the relevant Lie-algebra data. On the integrable-model side, $Q$-operators or transfer data play an analogous spectral role.

\paragraph{Target retained from the canonical source catalogue.}
\paragraph{Catalogue wording retained for traceability.}
What structural mechanism produces the QQ-tilde system for spectral determinants of affine opers, and is there an analogous system on the integrable-model side?

\paragraph{Source scope and qualification.}
The extraction does not include the functional relation, Lie type, shifts or normalizations. These are source-defined and must be supplied before using this record as a standalone equation; the two structural questions are retained without inventing a new QQ relation.
\subsection{Short English statement}
Explain the origin of the oper spectral-determinant functional relations and find their integrable-model counterpart.
\subsection{Sources}
\begin{itemize}
\item[\hypertarget{source:statement:1}{S1}] ULAM.AI and the UnsolvedMath Contributors, UnsolvedMath, catalogue entry. Catalogue curation under CC BY 4.0; definitions and English exposition added locally.
\url{https://huggingface.co/datasets/ulamai/UnsolvedMath}.
\item[\hypertarget{source:statement:2}{S2}] Algebraic Groups (2017), Oberwolfach Reports, DOI 10.4171/owr/2017/21, pp. 1281--1347. Primary report accessed 2026-10-11; relevant PDF pages 50.
\url{https://doi.org/10.4171/owr/2017/21}.
\item[\hypertarget{source:statement:3}{S3}] Frenkel et al., q-Opers and QQ-Systems (catalogue-listed context; not independently consulted in this preparation).
\url{https://arxiv.org/abs/2002.07344}.
\end{itemize}
\label{end:statement}
\end{document}
